\section{Marco Teórico}

\subsection{Señales de Voz: Producción y Características}
\label{sec:mt_voz}
Fundamentos fisiológicos y acústicos de la producción de voz humana. Representación digital de señales de audio (muestreo, cuantización, formato WAV). Parámetros relevantes: frecuencia de muestreo, duración, amplitud, energía.
 
\subsection{Análisis Espectral de Señales de Audio}
\label{sec:mt_espectral}
Transformada de Fourier (DFT/FFT). Transformada de Fourier de Tiempo Corto (STFT) y espectrograma. Escala Mel y banco de filtros Mel. Coeficientes Cepstrales de Frecuencia Mel (MFCC): cálculo e interpretación.

\subsection{Síntesis de Voz mediante Inteligencia Artificial}
\label{sec:mt_tts}
Sistemas Text-to-Speech (TTS) clásicos y neuronales. Arquitecturas relevantes: Tacotron 2, FastSpeech, VITS. Vocoders neuronales: WaveNet, HiFi-GAN, EnCodec (Meta). Características de la voz sintética que la diferencian de la natural.
 
\subsection{Aprendizaje Automático y Redes Neuronales}
\label{sec:mt_ml}
Conceptos fundamentales de aprendizaje supervisado. Redes neuronales artificiales: perceptrón multicapa (MLP). Funciones de activación, función de pérdida, retropropagación y descenso de gradiente.
 
\subsection{Redes Neuronales Profundas para Audio}
\label{sec:mt_dnn}
Redes Neuronales Convolucionales (CNN) aplicadas a espectrogramas. Redes Neuronales Recurrentes (RNN) y LSTM para secuencias temporales. Arquitecturas de atención y Transformers en procesamiento de audio. Modelos de aprendizaje auto-supervisado: Wav2Vec 2.0, HuBERT, WavLM.
 
\subsection{Detección de Voz Sintética (Anti-Spoofing)}
\label{sec:mt_antispoofing}
Definición del problema como clasificación binaria. Protocolo y métricas estándar: EER, t-DCF, AUC-ROC. Sistemas de referencia: GMM-LFCC, AASIST. Desafíos: generalización a vocoders desconocidos y degradaciones de canal.
 
\subsection{Corpus y Bases de Datos de Audio en Español}
\label{sec:mt_corpus}
Características y estructura de CIEMPIESS LIGHT, CORPUS CHM150 y TEDx Spanish Corpus. Criterios de selección y construcción de un corpus balanceado por género. Técnicas de segmentación de audio: bloques de duración fija.
 